\documentclass[11pt]{article}
\usepackage[utf8]{inputenc}
\usepackage{amsmath}
\usepackage{amsfonts}
\usepackage{amssymb}
\usepackage{bm}
\usepackage{geometry}
\usepackage{xcolor}
\usepackage{booktabs}
\usepackage{hyperref}

\geometry{margin=0.8in}

\title{Forward Spike Generation from a Reconstructed PRISM-EM Model\\
\large Ver3c aggregate and de-biased model replay}
\author{Algorithm Specification}
\date{\today}

\begin{document}

\maketitle

\begin{abstract}
This document specifies the planned
\texttt{forwardGenSpike\_ver3c.py} program. The program loads a fitted
PRISM-EM model produced either by the Stage (b) FDR-bag aggregator or the
Stage (c) de-biased fitter, replays its reconstructed latent-state timeline,
and draws a new spike train from the fitted switching Poisson GLM.
Stage (b) inputs use the Dale-pruned connectivity
\texttt{A\_prune} and fitted biases \texttt{B\_hat}; Stage (c) inputs use
\texttt{A\_debias} and \texttt{B\_debias}. The reconstructed state path is
continued cyclically, beginning at its final fitted state. The simulation
starts with an all-zero spike vector and preserves the exact log-rate
clipping threshold used during fitting. In addition to the generated spike
train, the program records the replayed states, the effective model, and a
time-resolved diagnostic counting how many neurons required log-rate
clipping at every generated step.
\end{abstract}

\tableofcontents
\newpage

% ===============================================================
\section{Purpose and Scope}
\label{sec:scope}
% ===============================================================

The forward generator tests the behavior implied by a reconstructed
network. It does not refit connectivity, rediscover states, or alter the
selected support. Instead, it treats the saved model as fixed and samples a
new realization of the fitted conditional point process.

For a requested time range \([t_L,t_R]\) in seconds, the program produces

\[
Y^{(f)} \in \mathbb{N}_0^{T_f\times N},
\]

where \(N\) is inferred from the fitted connectivity matrix. The bin width
\(\Delta t\), state-dependent bias matrix, reconstructed state timeline, and
upper log-rate clipping value are all inherited from the input fit. Thus the
number of generated bins and physical duration are

\[
T_f=\frac{t_R-t_L}{\Delta t}, \qquad
T_{\mathrm{seconds}}=t_R-t_L=T_f\,\Delta t.
\]
The range duration must contain an integer number of fitted time bins. The
simulation itself starts from the agreed zero-spike initial condition and
\texttt{S\_hat[-1]}; consequently, \(t_L\) names the requested range but
does not introduce a burn-in trajectory.

The generated files follow the version-2 NPZ schema implemented by
\texttt{toolbox/Util\_NumpyIOv2.py}. Schema-less legacy inputs and archives
whose declared shapes or dtypes disagree with their stored records are not
accepted.

% ===============================================================
\section{Command-Line Interface}
\label{sec:cli}
% ===============================================================

The planned command-line arguments are:
\begin{itemize}
  \item \texttt{--basePath}: analysis root containing
        \texttt{prismFit/} and receiving output under
        \texttt{spikesData/};
  \item \texttt{--inputModel}: input model stem, without
        \texttt{.prismEM.npz};
  \item \texttt{--time\_range\_sec START STOP}: requested output range in
        seconds. \texttt{STOP-START} must be positive and an integer
        multiple of the fitted \(\Delta t\);
  \item \texttt{--dataName}: optional output stem. Its default is
        \texttt{None};
  \item \texttt{--seed}: random seed for Poisson sampling, with a planned
        default of \texttt{42};
  \item \texttt{-v} or \texttt{--verbosity}: optional diagnostic
        verbosity.
\end{itemize}

The input path is

\[
\texttt{<basePath>/prismFit/<inputModel>.prismEM.npz}.
\]

If \texttt{--dataName} is supplied, it is used as the common output stem.
If it is omitted, the stem is

\[
\texttt{forwardN<N>\_<hash6>},
\]

where \(N\) is the dimension of the fitted \(A\) matrix and
\texttt{<hash6>} is a random six-character hexadecimal identifier. For
example, a model with \(N=200\) may produce
\texttt{forwardN200\_a1b2c3}.

An example invocation is:
\begin{verbatim}
basePath=/path/to/analysis

./forwardGenSpike_ver3c.py \
  --basePath $basePath \
  --inputModel daleN200_example_em1234_fdr1234_agr1234_debias \
  --time_range_sec 0 3600 \
  --seed 42
\end{verbatim}

% ===============================================================
\section{Accepted Input Models}
\label{sec:input}
% ===============================================================

The program accepts the two final-fit stages listed in
Table~\ref{tab:model-keys}. The fit type is determined from metadata and
confirmed by the required payload keys.

\begin{table}[ht]
\centering
\begin{tabular}{llll}
\toprule
Input stage & Metadata \texttt{fit\_type} & Connectivity & Biases \\
\midrule
Stage (b) aggregate & \texttt{prismEM\_FDRbags\_stageB}
  & \texttt{A\_prune} & \texttt{B\_hat} \\
Stage (c) de-biased & \texttt{prismEM\_deBias\_stageC}
  & \texttt{A\_debias} & \texttt{B\_debias} \\
\bottomrule
\end{tabular}
\caption{Stage-specific model parameters used for forward generation.}
\label{tab:model-keys}
\end{table}

Both input types must also provide:
\begin{itemize}
  \item \texttt{S\_hat}, the reconstructed hard state at every fitted
        time bin;
  \item \texttt{train.time\_step\_sec}, the fitted bin width
        \(\Delta t\);
  \item \texttt{train.eta\_clip}, the upper clipping threshold used by
        the fitted Poisson model;
  \item \texttt{train.num\_states} and
        \texttt{train.num\_neurons}, used to verify array shapes.
\end{itemize}

The top-level \texttt{poisson\_eta\_clip} metadata may be checked for
consistency, but \texttt{train.eta\_clip} is the authoritative value for
the model being replayed. The program must fail if the two available values
disagree materially, rather than silently changing the simulated model.

The following shape checks are required:
\[
A\in\mathbb{R}^{N\times N},\qquad
B\in\mathbb{R}^{M\times N},\qquad
\hat S\in\{0,\ldots,M-1\}^{T_{\mathrm{fit}}}.
\]
For a one-state model, a saved one-dimensional bias vector may be normalized
to shape \((1,N)\).

No explicit state-transition matrix is required. The saved model does not
currently contain a learned forward transition matrix, but this is not
missing information for the selected replay policy because the complete
reconstructed timeline \texttt{S\_hat} is present. Likewise, no source
spike vector is required because the requested initial spike state is zero.

% ===============================================================
\section{Reconstructed-State Replay}
\label{sec:state-replay}
% ===============================================================

Let the fitted hard state sequence be

\[
\hat S=(\hat S_0,\hat S_1,\ldots,\hat S_{L-1}),
\qquad L=T_{\mathrm{fit}}.
\]

The generated trajectory begins at the final reconstructed state and then
wraps around the saved timeline. For generated index
\(t=0,\ldots,T_f-1\),

\begin{equation}
\label{eq:state-wrap}
S^{(f)}_t = \hat S_{(L-1+t)\bmod L}.
\end{equation}

Therefore \(S^{(f)}_0=\hat S_{L-1}\), the next bin uses \(\hat S_0\), and
the pattern repeats as many times as needed. If the number of bins derived
from \texttt{--time\_range\_sec} is shorter than \(L\), only the requested
prefix of this cyclic continuation is used. If it exceeds \(L\), the same
state timeline is reused by modular indexing without allocating multiple
full copies in advance.

The Stage (b) locked M-step and the usual locked Stage (c) refit train the
bias matrix against hard, one-hot destination-bin state labels. Forward
generation therefore uses

\[
C^{(f)}_{tm}=\mathbf{1}\{S^{(f)}_t=m\}
\]

and does not mix bias atoms using the continuous \texttt{c\_hat} values.
The one-hot coefficient matrix is saved for traceability in the companion
forward-truth file.

% ===============================================================
\section{Forward Poisson Model}
\label{sec:model}
% ===============================================================

The simulation begins from an all-zero spike history,

\[
Y^{(f)}_{-1}=\bm{0}.
\]

At each generated time step, compute the raw log-rate vector

\begin{equation}
\label{eq:raw-eta}
\eta^{\mathrm{raw}}_t
= A Y^{(f)}_{t-1}+B_{S^{(f)}_t},
\end{equation}

or, equivalently for row-vector implementations,

\[
\eta^{\mathrm{raw}}_t
=Y^{(f)}_{t-1}A^{\mathsf T}+C^{(f)}_tB.
\]

The clipping operation must match fitting exactly: only the upper tail is
clipped,

\begin{equation}
\label{eq:eta-clip}
\eta_{ti}=\min\left(\eta^{\mathrm{raw}}_{ti},
\eta_{\mathrm{clip}}\right).
\end{equation}

The conditional mean spike count and sampled count are

\begin{equation}
\label{eq:poisson-forward}
\lambda_{ti}=\exp(\eta_{ti})\,\Delta t,
\qquad
Y^{(f)}_{ti}\sim\operatorname{Poisson}(\lambda_{ti}).
\end{equation}

The same seeded NumPy random generator is used for every step. The state
timeline is deterministic for a fixed input model and requested length;
\texttt{--seed} controls only the stochastic Poisson realization. Generated
counts are stored as \texttt{int32}, avoiding the \texttt{uint8} saturation
used by some compact experimental preprocessing files.

% ===============================================================
\section{Tracking Log-Rate Clipping}
\label{sec:clipping}
% ===============================================================

Clipping must be measured before Eq.~\eqref{eq:eta-clip} is applied. Define

\begin{equation}
\label{eq:clip-count}
K_t=\sum_{i=1}^{N}
\mathbf{1}\{\eta^{\mathrm{raw}}_{ti}>\eta_{\mathrm{clip}}\}.
\end{equation}

The integer vector \texttt{num\_clipped\_neurons\_time} stores \(K_t\) for
all \(T_f\) generated bins. It has range \(0\le K_t\le N\). Equality with
the threshold is not counted as clipping because the value is unchanged by
the upper clamp.

The exact histogram is

\begin{equation}
\label{eq:clip-hist}
H_k=\sum_{t=0}^{T_f-1}\mathbf{1}\{K_t=k\},
\qquad k=0,1,\ldots,N.
\end{equation}

The forward-truth file stores:
\begin{itemize}
  \item \texttt{num\_clipped\_neurons\_time}: \(K_t\), shape \((T_f,)\);
  \item \texttt{clipped\_neuron\_hist}: \(H_k\), shape \((N+1,)\);
  \item \texttt{clipped\_neuron\_hist\_fraction}:
        \(H_k/T_f\), shape \((N+1,)\);
  \item \texttt{clipped\_neuron\_hist\_k}: integer support
        \(0,1,\ldots,N\);
  \item \texttt{clipped\_neuron\_quartiles}: the four upper quartile
        cut points
        \((q_{25},q_{50},q_{75},q_{100})\) of the \(K_t\) distribution.
\end{itemize}

Here

\[
q_p=\operatorname{quantile}\left(\{K_t\}_{t=0}^{T_f-1},p/100\right),
\qquad p\in\{25,50,75,100\}.
\]

The histogram and the four quartiles are printed at the end of the run.
Because \(K_t\) is integer-valued, interpolated quartile values may be
non-integer; the saved histogram remains the authoritative exact summary.
The summary should also print the number and fraction of time bins with
any clipping, making sustained saturation immediately visible.

% ===============================================================
\section{Output Files}
\label{sec:outputs}
% ===============================================================

Let \texttt{<dataName>} denote the supplied or automatically generated core
name. The program writes:

\[
\texttt{<basePath>/spikesData/<dataName>.spikes.npz},
\]
\[
\texttt{<basePath>/spikesData/<dataName>.forwardTruth.npz}.
\]

\subsection{Generated spike file}

The \texttt{.spikes.npz} payload contains:
\begin{itemize}
  \item \texttt{spikes}: generated counts, shape \((T_f,N)\),
        \texttt{int32};
  \item \texttt{single\_rates}: generated mean rate per neuron in Hz,
        computed as \(\operatorname{mean}_t(Y^{(f)}_{ti})/\Delta t\).
\end{itemize}

Its metadata records at least the data type, time step, fit-time
\texttt{poisson\_eta\_clip}, neuron count, random seed, output name, input
model name and file, input fit type, selected \(A/B\) keys, and
\texttt{num\_steps}. Provenance must be copied rather than mutated in place
and augmented with the forward-generation program and model source.

\subsection{Forward-truth file}

The companion \texttt{.forwardTruth.npz} payload contains:
\begin{itemize}
  \item \texttt{A\_forward} and \texttt{B\_forward}: exact arrays used by
        Eqs.~\eqref{eq:raw-eta}--\eqref{eq:poisson-forward};
  \item \texttt{S\_forward}: wrapped hard-state sequence, shape
        \((T_f,)\);
  \item \texttt{C\_forward}: corresponding one-hot states, shape
        \((T_f,M)\);
  \item the clipping timeline, histogram, histogram support, normalized
        histogram, and quartiles defined in Section~\ref{sec:clipping}.
\end{itemize}

Its metadata records the state replay rule
\texttt{last\_state\_then\_cyclic\_wrap}, the zero-spike initial condition,
the input state-history length, the number of complete wraps, the remainder
length, \(N\), \(M\), \(\Delta t\), \(\eta_{\mathrm{clip}}\), seed, requested
time range and duration, and stage-specific parameter selection.

Both archives are written through \texttt{Util\_NumpyIOv2} and therefore
contain \texttt{schema.JSON}. Metadata are normalized with
\texttt{json\_safe\_metadata} before writing.

\subsection{Visualization}

The generated pair can be displayed with
\texttt{view\_spikesTrain3.py}; a typical invocation is:
\begin{verbatim}
./view_spikesTrain3.py --basePath $basePath --dataName <dataName> \
  --idxState -1 --time_range_sec 0 15 -p b
\end{verbatim}
The viewer detects \texttt{data\_type=forwardPrism}, loads the matching
\texttt{.forwardTruth.npz} archive, and displays the generated activity,
wrapped reconstructed state, and number of eta-clipped neurons at every
time step. Ordinary multi-state inputs continue to use their
\texttt{.prismTruth.npz} companion and oracle-state panel.

% ===============================================================
\section{Algorithm Summary}
\label{sec:summary}
% ===============================================================

The complete forward workflow is:
\begin{enumerate}
  \item Resolve and schema-validate the input model archive.
  \item Identify Stage (b) or Stage (c) and select the corresponding
        \(A/B\) arrays from Table~\ref{tab:model-keys}.
  \item Validate \(A\), \(B\), \texttt{S\_hat}, \(N\), \(M\),
        \(\Delta t\), and \(\eta_{\mathrm{clip}}\).
  \item Choose the output stem from \texttt{--dataName}, or construct
        \texttt{forwardN<N>\_<hash6>}.
  \item Build \(S^{(f)}\) by Eq.~\eqref{eq:state-wrap}, beginning at
        \texttt{S\_hat[-1]} and wrapping as needed; form one-hot
        \(C^{(f)}\).
  \item Initialize the previous spike vector to zero.
  \item For every output bin, evaluate the raw log-rate, count clipped
        neurons, apply the inherited upper clip, sample Poisson counts, and
        use those counts as the next lagged input.
  \item Compute generated single-neuron rates, the clipping histogram, its
        normalized form, the four quartiles, and the fraction of bins with
        clipping.
  \item Write the schema-v2 spike and forward-truth archives and print their
        paths and clipping summaries.
\end{enumerate}

No model fitting occurs in this workflow. Reproducibility requires an
unchanged input archive and the same \texttt{--time\_range\_sec} and
\texttt{--seed}; the random output hash affects only file naming.

% ===============================================================
\section{Required Failure Checks}
\label{sec:checks}
% ===============================================================

The program should stop with an explicit error if:
\begin{itemize}
  \item the input archive is missing or is not schema version 2;
  \item its fit type is neither a Stage (b) aggregate nor Stage (c)
        de-biased fit;
  \item the required stage-specific \(A/B\) arrays are absent;
  \item \texttt{S\_hat} is absent, empty, non-integral, or contains a state
        index outside \([0,M-1]\);
  \item \(A\), \(B\), or metadata dimensions disagree;
  \item \texttt{--time\_range\_sec} is non-finite, has a negative start,
        does not satisfy \texttt{STOP > START}, or its duration does not
        contain an integer number of fitted time bins;
  \item \(\Delta t\) is not positive, or the fit-time clipping threshold is
        missing or non-finite;
  \item an output path already exists, unless an explicit overwrite policy
        is added later.
\end{itemize}

These checks ensure that forward generation cannot silently substitute a
different connectivity array, bias matrix, time scale, clipping threshold,
or state convention.

\end{document}
